% Lösungen Einheit 1 von 3 – Vektorbegriff, Betrag und Kollinearität (zusammenbau v0.1)
\einheitenkopf{Einheit 1 von 3 · Vektorbegriff, Betrag und Kollinearität}

\erg{9}{a) Zahl – eine Koordinate; die Länge liefert sie.}
\erg{10}{a) Zahl \quad b) $\overrightarrow{AB} = (3 | 1 | 2)$, $\overrightarrow{BA} = (-3 | -1 | -2)$, $2 \cdot \overrightarrow{AB} = (6 | 2 | 4)$ \quad c) $|\vec u| = 7$, $|\vec v| = 6$; $\vec u$ ist länger. \quad d) $a^2 + 144 = 169$, $a^2 = 25$, wegen $a > 0$ ist $a = 5$ \quad e) $\overrightarrow{DC} = k \cdot \overrightarrow{AB}$ mit $\overrightarrow{AB} = (0 | 3 | 2)$; aus der dritten Koordinate $4 = 2k$, $k = 2$; dann $y - 1 = 6$, also $y = 7$ \quad f) z. B. $(1 | 0 | 0)$; von $(0 | 1 | 0)$ aus sieht man nur eine Strecke. Von oben: ein X aus den Strecken von $(0 | 0)$ nach $(4 | 4)$ und von $(0 | 4)$ nach $(4 | 0)$. \quad g) Ja: $(600 | 0 | 0)$ hat den Betrag 600, $(200 | 200 | 200)$ den Betrag $\sqrt{120\,000} \approx 346{,}4$ – gleiche Summe, verschiedene Längen.}

10f)

\begin{ksys}[xmin=0,xmax=5,ymin=0,ymax=5,klein] \funktionab{\x}{}{0}{4} \funktionab{4-\x}{}{0}{4} \end{ksys}

\erg{11}{a) Fehler: die Koordinaten addiert statt quadriert und gewurzelt. Richtig: $\sqrt{16 + 16 + 49} = \sqrt{81} = 9$. \quad b) Weil $5^2$ und $(-5)^2$ beide 25 ergeben – das Quadrieren löscht das Vorzeichen. \quad c) $\overrightarrow{AB} = (3 | 2)$ \quad d) $|\overrightarrow{SZ}| = \sqrt{900 + 1600} = 50$ m; nein, 45 m reichen nicht.}
